% META: {"id": "TSPE-CONTIN-CO-038", "chapitre": "TSPE-CONTINUITE", "type_objet": "corrige", "exercice_ref": "TSPE-CONTIN-EX-038", "capacites_codes": ["C2"], "sources_inspiration": [], "status": "generated"}
% BEGIN-VERIFY
% from sympy import symbols, diff, Rational, sqrt, simplify, solveset, S as SS, Eq, nsolve
% x = symbols('x')
% def iterer(f, u0, n):
%     u = u0
%     out = [u]
%     for _ in range(n):
%         u = f.subs(x, u)
%         out.append(u)
%     return out
% f = sqrt(2*x)
% assert solveset(Eq(f, x), x, SS.Reals) == {0, 2}
% t = [float(v) for v in iterer(f, 1, 3)]
% assert abs(t[1] - 1.4142136) < 1e-6
% assert abs(t[2] - 1.6817928) < 1e-6
% assert abs(t[3] - 1.8340081) < 1e-6
% assert all(1 <= float(f.subs(x, v)) <= 2 for v in (1, 1.5, 2))
% assert float(f.subs(x, 1)) == 2**0.5 and float(f.subs(x, 2)) == 2.0
% assert all(float((f - x).subs(x, v)) > 0 for v in (1, 1.5, 1.9))
% assert simplify((f - x).subs(x, 2)) == 0
% # la suite issue de u0 = 0 est constante nulle
% assert [float(v) for v in iterer(f, 0, 3)] == [0.0, 0.0, 0.0, 0.0]
% END-VERIFY
\begin{corrige}{TSPE-CONTIN-EX-038}

\textbf{Q1.} $u_1 = \sqrt2 \approx 1{,}414$ ; $u_2 = \sqrt{2\sqrt2} \approx 1{,}682$ ; $u_3 \approx 1{,}834$.

\textbf{Q2.} Sur $[0\,;+\infty[$, $\sqrt{2x} = x$ equivaut, en elevant au carre (les deux membres etant positifs), a $2x = x^2$, soit $x(x-2) = 0$. Les points fixes sont donc $0$ et $2$ : il y en a \textbf{deux}.

\textbf{Q3.} Si $x \in [1\,;2]$ alors $2 \leq 2x \leq 4$, donc $\sqrt2 \leq \sqrt{2x} \leq 2$, et $f(x) \in [1\,;2]$ puisque $\sqrt2 > 1$.

Recurrence : $u_0 = 1 \in [1\,;2]$, et la stabilite propage l'encadrement a tout rang.

\textbf{Q4.} Sur $[1\,;2]$, $f(x)$ et $x$ sont positifs, donc $f(x) - x$ a le signe de $2x - x^2 = -x(x-2)$, strictement positif sur $[1\,;2[$ et nul en $2$. La suite est donc croissante ; majoree par $2$, elle converge vers $\ell$. Par continuite de $f$ sur $[1\,;2]$, $\ell$ verifie $f(\ell) = \ell$, donc $\ell \in \{0\,;2\}$ d'apres Q2. Comme $u_n \geq 1$ pour tout $n$, la limite verifie $\ell \geq 1$ : on conclut $\ell = 2$.

\textbf{Q5.} La resolution de $f(\ell) = \ell$ ne selectionne pas la limite : elle fournit seulement l'ensemble des \emph{candidats}. C'est l'encadrement $u_n \geq 1$, obtenu independamment, qui elimine $0$. Autrement dit, la continuite donne une condition necessaire, pas suffisante.

Une suite du meme type convergeant vers $0$ est la suite constante $u_0 = 0$, puisque $f(0) = 0$ : elle demarre sur le point fixe et n'en bouge plus. Toute suite demarrant en $u_0 > 0$ est attiree vers $2$ ; le point fixe $0$ est dit \emph{repulsif}.
\end{corrige}
